\clearpage
\subsection{\titlepkg{thm} 模块}
本模块主要用于定理类以及证明类数学环境定制. 本模块提供了丰富的接口以及选项，与此同时本模块提供了丰富的
Hook, 方便用户直接对环境进行操作.

\pkg{thm} 提供的数学环境主要分为两类:
\begin{itemize}
  \item 定理类: \env{axiom, definition, theorem, lemma, corollary, proposition, remark};
  \item 证明类: \env{proof, exercise, example, solution, problem}
\end{itemize}

所以请区分 ``定理类'' 和 ``证明类'' 两类环境, 以便于正确地使用 \pkg{thm} 提供的各个命令. \ztex{} 的 \pkg{thm} module 
中的部分命令或变量也许没有显式地含有 \textcolor{red}{\sffamily theorem} 字样, 但是这些命令或变量仍然是属于 ``定理类'' 的; 应用于
``证明类'' 环境的命令或变量均显式地含有 \textcolor{red}{\sffamily proof} 字样.


\clearpage
\subsubsection{用户接口}
\begin{function}[updated=2024-11-05]{\qedsymbol}
  \begin{syntax}
    \cs{qedsymbol}
  \end{syntax}
  此命令用于输出证明环境的结束符号, 默认为 $\square$.
\end{function}


\begin{function}[updated=2025-04-25]{\zthmlang}
  \begin{syntax}
    \cs{zthmlang}\marg{lang}
  \end{syntax}
  此命令用于设置定理类环境的语言(从而会影响到其标题名称), 目前支持 \texttt{cn, en, fr} 三种语言, 仅能在文档的导言区使用.\par
  一个使用样例如下\zchcmd:
\end{function}
\vskip1.25em\relax
\begingroup
\begin{DocExample}*
  \begin{theorem}[zthmlang-1]
    This is a chinese zthmlang-1.
  \end{theorem}
  \zthmlang{fr}
  \begin{theorem}[zthmlang-2]
    This is a france zthmlang-2.
  \end{theorem}
  \zthmlang{en}
  \begin{theorem}[zthmlang-3]
    This is a english zthmlang-3.
  \end{theorem}
\end{DocExample}
\endgroup


\begin{function}[updated=2025-04-25]{\zthmnameset}
  \begin{syntax}
    \cs{zthmnameset}\marg{lang}\marg{key-value}
  \end{syntax}
  此命令用于设置数学环境的名称, 包括 ``定理类'' 和 ``证明类'', 仅能在文档的导言区使用. 预定义的 \meta{lang} 值有: \texttt{en, cn, fr}. 除预定义的
  这三种语言外, 用户可以使用此命令自行声明(\meta{lang}), 然后使用命令 \cs{zthmlang}\texttt{\{\meta{lang}\}} 进行切换.% 
  \textbf{注意}:此命令需应用于 \cs{zthmlang} 命令之前, 否则此命令的相关设置将不会生效; 此命令无法控制由 \cmd{\zthmenvnew} 
  或 \cmd{\zthmenvset} 创建的环境标题.
\end{function} 

\vspace{2em}
\noindent\fbox{\parbox{1\linewidth}{下面我们采用键值对的方式对 \meta{key-value} 这一项参数进行描述: \texttt{zthmnameset/} 
表示它是此 \meta{key-value} 参数的父级命令; 后续为了行文的方便，我们在描述一个(父级)命令之后，使用 \texttt{../} 
来表示其缩写形式(\texttt{../} 有时也用于表示任意的键名, 即由用户定义的键名).\par
\textbf{注意}:虽然它的设置方法和 \texttt{key-value} 这样的数据结构类似, 但是用户不能将 \cs{keys_define:nn} 这样的
命令应用于这类键值对, 而应使用其父级命令 \cs{zthmnameset} 对其进行设置.}}

\begin{keyval}[parent=zthmnameset]{
    axiom, definition, theorem, lemma, corollary, proposition, remark,
    proof, exercise, example, solution, problem
  }
  \begin{syntax}
    axiom       = \marg{名称}>\dval{Axiom}
    definition  = \marg{名称}>\dval{Definition}
    theorem     = \marg{名称}>\dval{Theorem}
    lemma       = \marg{名称}>\dval{Lemma}
    corollary   = \marg{名称}>\dval{Corollary}
    proposition = \marg{名称}>\dval{Proposition}
    remark      = \marg{名称}>\dval{Remark}
    proof       = \marg{名称}>\dval{Proof}
    exercise    = \marg{名称}>\dval{Exercise}
    example     = \marg{名称}>\dval{Example}
    solution    = \marg{名称}>\dval{Solution}
    problem     = \marg{名称}>\dval{Problem}
  \end{syntax}
  当 \meta{lang}\texttt{ = en} 时, 命令 \cs{zthmnameset} 中 \meta{key-value} 的设置情况.
\end{keyval}

\begin{keyval}[parent=..]{
    axiom, definition, theorem, lemma, corollary, proposition, remark,
    proof, exercise, example, solution, problem
  }
  \begin{syntax}
    axiom       = \marg{名称}>\dval{Axiome}
    definition  = \marg{名称}>\dval{Définition}
    theorem     = \marg{名称}>\dval{Théorème}
    lemma       = \marg{名称}>\dval{Lemme}
    corollary   = \marg{名称}>\dval{Corollaire}
    proposition = \marg{名称}>\dval{Proposition}
    remark      = \marg{名称}>\dval{Remarque}
    proof       = \marg{名称}>\dval{Preuve}
    exercise    = \marg{名称}>\dval{Exercice}
    example     = \marg{名称}>\dval{Exemple}
    solution    = \marg{名称}>\dval{Solution}
    problem     = \marg{名称}>\dval{Problème}
  \end{syntax}
  当 \meta{lang}\texttt{ = fr} 时, 命令 \cs{zthmnameset} 中 \meta{key-value} 的设置情况.
\end{keyval}


\begin{keyval}[parent=..]{
    axiom, definition, theorem, lemma, corollary, proposition, remark,
    proof, exercise, example, solution, problem
  }
  \begin{syntax}
    axiom       = \marg{名称}>\dval{公理}
    definition  = \marg{名称}>\dval{定义}
    theorem     = \marg{名称}>\dval{定理}
    lemma       = \marg{名称}>\dval{引理}
    corollary   = \marg{名称}>\dval{推论}
    proposition = \marg{名称}>\dval{命题}
    remark      = \marg{名称}>\dval{备注}
    proof       = \marg{名称}>\dval{证明}
    exercise    = \marg{名称}>\dval{练习}
    example     = \marg{名称}>\dval{示例}
    solution    = \marg{名称}>\dval{解}
    problem     = \marg{名称}>\dval{问题}
  \end{syntax}
  当 \meta{lang}\texttt{ = cn} 时, 命令 \cs{zthmnameset} 中 \meta{key-value} 的设置情况.
\end{keyval}
一个基本的使用案例如下\zchcmd:
\begin{DocExample}*
  \zthmnameset{cn}{
    theorem=新定理,
    proof=证
  }
  \zthmlang{cn} % update thm title name
  \begin{theorem}[zthmnameset-1]
    This is a theorem zthmnameset-1.
  \end{theorem}
  \begin{proof}
    This is a proof.
  \end{proof}
\end{DocExample}


\zthmstyle{background}
\begin{function}[added=2025-09-23]{\zthmenvnew}
  \begin{syntax}
    \cs{zthmenvnew}\oarg{type}\marg{key-value}
  \end{syntax}
  根据第二个参数中的 \meta{key-value} 创建一系列类型为 \meta{type} 的定理环境, 仅可在导言区使用; 
  如果对应的环境已存在，则覆盖其原有的定义. \meta{type} 可选 \texttt{theorem, proof} 两种类型, 默认
  为 \texttt{theorem}. 键值 \meta{key-value} 的语法请参见下面的样例:
  \textcolor{red}{\bfseries 例 \inteval{\theDocExample+1}}.
  % \textbf{注意}: 上述格式中的 `\texttt{\string|}' 不可以省略, 否则会导致解析错误.\par
  % 一个基本的使用案例如下\zchcmd:
\end{function}


\begin{keyval}[parent=ztex/thmnew]{name, color, tocsym}
  \begin{syntax}
    name   = \meta{title} > \dval{\meta{name}}
    color  = \meta{color} > \dval{black}
    tocsym = \meta{code} > \dval{\meta{name}}
  \end{syntax}
  \meta{title} 为新环境对应的名称, 可以省略, 默认以此环境的名称为标题; 
  \meta{color} 为合法的 \ztex{} 色彩格式, 可以省略, 默认为 ``\texttt{black}'';
  \meta{tocsym} 为该环境在定理目录中的前缀, 若为空或未设置此键, 则默认为此环境的名称.
  \textbf{备注}: 用户也可以使用 \cmd{\zthmtocsym} 命令来设置新环境的 \meta{tocsym}.
\end{keyval}


\begin{function}[added=2025-09-23]{\zthmenvset}
  \begin{syntax}
    \cmd{\zthmenvset}\Arg{type}\Arg{key-value}
  \end{syntax}
  此命令和上述的 \cmd{\zthmenvnew} 类似, 但此命令会覆盖原有的环境定义, 且可以在导言区或正文中使用.
\end{function}


\begin{DocExample}*
\zthmenvnew{
  Zaxiom, 
  Ztheorem={name=Thm, color={HTML}{a0d911}, tocsym=\textbf{ZT}\;},
  Zproposition={name=Prop, color=blue},
}
\zthmenvnew[proof]{
  Zproof={color={cyan}},
  Zexample={name=EXAMPLE, color=red},
  Zsolution={name=Solution},
}
\begin{Zproof}[zthmenvnew-1]
  This is a Zproof zthmenvnew-1.
\end{Zproof}
\begin{Zexample}[zthmenvnew-2]
  This is a Zexample zthmenvnew-2.
\end{Zexample}
\begin{Ztheorem}[zthmenvnew-3]
  This is a Ztheorem zthmenvnew-3
\end{Ztheorem}
\end{DocExample}
\zthmstyle{plain}


\begin{function}[updated=2025-04-25]{\zthmcnt}
  \begin{syntax}
    \cs{zthmcnt}\marg{key-value}
  \end{syntax}
  此命令用于定义数学类环境的计数器, 仅能在导言区使用.
\end{function}


\begin{keyval}[parent=..]{parent, share}
  \begin{syntax}
    parent = \meta{counter}>\dval{section}
    share  = \meta{true|\textbf{false}}>\dval{false}
  \end{syntax}
  \meta{parent} 用于指定定理类环境计数器的父计数器, 默认父计数器为 \texttt{section}; 当父计数器更新时，此环境
  的计数器便会重置; \meta{share} 用于控制所有的定理类环境是否共用一个计数器，默认为 \texttt{false}. \textbf{注意}:%
  若指定所有定理类环境公用计数器，此时 \cs{cref} 对应的共同名称为 ``result'' 或 ``结果'', 具体取决于 \cs{zthmlang} 的设置.
\end{keyval}


\begin{function}[updated=2025-04-25]{\zthmstyle}
  \begin{syntax}
    \cs{zthmstyle}\marg{style}
  \end{syntax}
  此命令用于设置定理类环境的样式, 仅能在导言区使用. \textbf{注意}:\textcolor{red}{\sffamily 由于技术原因, 当用户需要加载
  \pkg{thm} library 时, 必须将命令 \cs{zthmstyle}\texttt{\marg{style}} 置于 \cs{ztexloadlib}\zarg{thm} 之前}. 
\end{function}


\begin{keyval}[parent=ztex/thm/style]{plain, leftbar, background, fancy}
  \begin{syntax}
    plain > \nval
    leftbar > \nval
    background > \nval
    fancy > \nval
  \end{syntax}
\end{keyval}
一个基本的使用样例如下\zchcmd:
\begin{DocExample}*
  \zthmstyle{plain}
  \begin{theorem}[zthmstyle-1]
    A `plain' style zthmstyle-1.
  \end{theorem}
  \zthmstyle{leftbar}
  \begin{theorem}[zthmstyle-2]
    A `leftbar' style zthmstyle-2.
  \end{theorem}
  \zthmstyle{background}
  \begin{theorem}[zthmstyle-3]
    A `background' style zthmstyle-3.
  \end{theorem}
  \zthmstyle{fancy}
  \begin{theorem}[zthmstyle-4]
    A `fancy' style zthmstyle-4.
  \end{theorem}
\end{DocExample}
\zthmstyle{plain}


\begin{function}[updated=2025-04-25]{\zthmcolorset}
  \begin{syntax}
    \cs{zthmcolorset}\marg{key-value}
  \end{syntax}
  此命令和 \cmd{\zcolorset} 类似，但其仅用于对数学环境的色彩设置(比如, 你不能在此命令中设置 \meta{link} 对应的色彩)，
  且仅能在导言区使用. 此命令仅能用于数学类环境的色彩自定义, 如果出现除数学(包括由命令 \cs{zthmnew} 所创建的)环境以外色彩设置，
  那么 \ztex{} 会抛出错误;
\end{function}


\begin{keyval}[parent=..]{
    axiom, definition, theorem, lemma, corollary, proposition, remark,
    proof, exercise, example, solution, problem
  }
  \begin{syntax}
    axiom       = \meta{color spec}>\dval{ztex@color@axiom}
    definition  = \meta{color spec}>\dval{ztex@color@definition}
    theorem     = \meta{color spec}>\dval{ztex@color@theorem}
    lemma       = \meta{color spec}>\dval{ztex@color@lemma}
    corollary   = \meta{color spec}>\dval{ztex@color@corollary}
    proposition = \meta{color spec}>\dval{ztex@color@proposition}
    remark      = \meta{color spec}>\dval{ztex@color@remark}
    ...
  \end{syntax}
  \meta{color spec} 为一个合法的 \ztex{} 色彩格式. ``\texttt{proof, exercise, ..., problem}''
  的默认值均为 ``\texttt{black}''.
\end{keyval}



\clearpage
\subsubsection{定理目录}
本节的命令依赖于 \ztex{} 的 \pkg{zsect} 模块, 欲使用本节的命令, 请务必加载 \pkg{sect} 模块. 在加载 \pkg{sect} 模块后, 定理目录样式的定制请参考 \pkg{sect} 模块的文档.

\begin{function}[updated=2026-02-12]{\zthmtoc}
  \begin{syntax}
    \cs{zthmtoc}\oarg{keyval}
  \end{syntax}
  此命令用于打印定理类环境对应的目录, 可选参数 \meta{keyval} 用于自定义定理目录的样式, 配置方法请参见 \cref{ss2:toc} 的文档说明, 默认为空. \textbf{备注}: 定理目录的默认层级为 ``\texttt{subsection}'', 证明类环境不会计入定理目录; 此命令依赖于 \pkg{sect} 模块. 该命令的一个简单使用样例如下:
\end{function}
\begin{DocExample}*
\zthmtoc[
  page.format+ = \texttt,
  hyper.title  = true,
]
\begin{proposition}[zthmtoc-1]proposition zthmtoc-1\end{proposition}
\begin{lemma}[zthmtoc-2]lemma zthmtoc-2\end{lemma}
\begin{corollary}[zthmtoc-3]corollary zthmtoc-3\end{corollary}
\end{DocExample}


\begin{function}[updated=2025-04-25]{\zthmtocadd}
  \begin{syntax}
    \cs{zthmtocadd}\oarg{class}\marg{key-value}
  \end{syntax}
  此命令用于向定理类环境目录中添加条目, \meta{class} 表示该条目在目录中的层级,  默认为 ``\texttt{section}''. 其它的可选值有: \texttt{chapter}, \texttt{subsection} 等.
  \textbf{备注}: 此命令依赖于 \pkg{sect} 模块.
\end{function}

\begin{keyval}[parent=..]{name}
  \begin{syntax}
    name = \marg{条目名称}>\dval{无}
  \end{syntax}
  目前的键仅有 \texttt{name}, 后续可能有变动.
\end{keyval}
一个简单的使用样例如下:
\begin{DocExample}*
\zthmtocadd[section]{name=New:Added Thm ITEM}
\end{DocExample}



\begin{function}[updated=2025-09-04]{\zthmtocstop}
  \begin{syntax}
    \cs{zthmtocstop}
  \end{syntax}
  此命令用于停止向定理类环境目录中添加条目, 用户也可以使用 \cmd{\ztocstop}\texttt{[lom]} 替代此命令. \textbf{备注}: 此命令依赖于 \pkg{sect} 模块.
\end{function}



\begin{function}[updated=2025-04-25]{\zthmtoclevel}
  \begin{syntax}
    \cs{zthmtoclevel}\marg{class}
  \end{syntax}
  此命令用于设置定理目录中各条目的层级, 仅能在导言区使用, \meta{class} 可以为 ``\texttt{chapter, section, subsection, ...}'' 等. \textbf{备注}: 目前暂不支持将不同的定理类环境名设置为不同的层级.
\end{function}


\begin{function}[updated=2025-04-25]{\zthmtocprefix}
  \begin{syntax}
    \cs{zthmtocprefix}\marg{prefix}
  \end{syntax}
  此命令用于所有定理类环境目录中所有条目的共同前缀, 默认为空. 
\end{function}


\begin{function}[updated=2025-04-25]{\zthmtocsym}
  \begin{syntax}
    \cs{zthmtocsym}\marg{key-value}
  \end{syntax}
  此命令用于分别设置所有定理类环境名在目录中的前缀, 仅能在导言区使用.
\end{function}

\begin{keyval}[parent=..]{axiom, definition, theorem, lemma, corollary, proposition, remark}
  \begin{syntax}
    axiom       = \meta{前缀}>\dval{A\char92\vsp }
    definition  = \meta{前缀}>\dval{D\char92\vsp }
    theorem     = \meta{前缀}>\dval{T\char92\vsp }
    lemma       = \meta{前缀}>\dval{L\char92\vsp }
    corollary   = \meta{前缀}>\dval{C\char92\vsp }
    proposition = \meta{前缀}>\dval{P\char92\vsp }
    remark      = \meta{前缀}>\dval{R\char92\vsp }
  \end{syntax}
  其中 \meta{前缀} 为任意合法的 \LaTeX{} 代码.
\end{keyval}
此命令的大致使用方法如下:
\begin{DocExample}
\zthmtocsym{
  axiom        =  AA,
  definition   =  DD,
  theorem      =  TT,
  lemma        =  LL,
  corollary    =  CC,
  proposition  =  PP,
  remark       =  RR,
}
\end{DocExample}


\begin{function}[updated=2025-04-25]{\zthmtocsymrm}
  此命令用于清除所有由命令 \cmd{\zthmtocsym} 添加在目录中的前缀, 仅能在导言区使用. 
  \textbf{注意}:此命令不能清除由 \cmd{\zthmtocprefix} 指定的前缀.
\end{function}



\clearpage
\subsubsection{高级接口}
\begin{function}[added=2025-09-03]{\zthmremoveCJKecglue, \zthmrestoreCJKecglue}
  命令 \cmd{\zthmremoveCJKecglue} 用于取消原始的 CJKecglue,
  命令 \cmd{\zthmrestoreCJKecglue} 用于恢复原始的 CJKecglue 设置.
\end{function}

\begin{function}[EXP, updated=2024-11-05]{\zthmnumber}
  此命令表示对应环境的编号, 类似于 \pkg{amsthm} 中的 \cmd{\thmnumber}. 用户不应在除 \cmd{\zthmtitleformat} 外的
  任何地方使用, 在命令 \cs{zthmtitleformat} 之外, 此命令输出的内容无任何实际意义.
\end{function}

\begin{function}[EXP, updated=2024-11-05]{\zthmname}
  此命令表示对应环境的名称, 类似于 \pkg{amsthm} 中的 \cmd{\thmname}. 用户不应在除 \cmd{\zthmtitleformat} 外的
  任何地方使用, 在命令 \cs{zthmtitleformat} 之外, 此命令输出的内容无任何实际意义.
\end{function}

\begin{function}[EXP, updated=2024-12-05]{\zthmnote}
  \begin{syntax}
    \cs{zthmnote}\marg{prefix}\marg{suffix}
  \end{syntax}
  此命令表示对应环境的注释, 类似于 \pkg{amsthm} 中的 \cmd{\thmnote}. 用户不应在除 \cmd{\zthmtitleformat} 外的
  任何地方使用, 在命令 \cs{zthmtitleformat} 之外, 此命令输出的内容无任何实际意义.
\end{function}


\begin{function}[updated=2025-04-25, EXP]{\thm@tmp@name}
  此命令用于临时保存定理类环境的名称, 用户可以在自定义定理类环境样式时使用. \textbf{注意}: 此命令和前述的
  \cmd{\zthmname} 不同, 因 \cs{thm@tmp@name} 只能取值于合法的定理类环境名称集合, 而 \cmd{\zthmname}
  是 \cs{thm@tmp@name} 的格式化版本, 可能包含 \cs{bfseries}, \cs{sffamily} 等格式化命令.
\end{function}


\begin{function}[updated=2025-04-25, EXP]{\thm@tmp@color, \thmproof@tmp@color}
  此二命令用于临时保存定理类环境和证明类环境的色彩, 用于在 \cmd{\zthmtitleformat} 中进行色彩切换.%
  \textbf{注意}:普通用户在使用这两个命令时, 请将其置于 \cs{makeatletter} 和 \cs{makeatother} 之间.
\end{function}


\begin{function}[EXP, updated=2024-11-05]{\zthmtitle, \zthmtitle*}
  \cs{zthmtitle} 命令为定理类环境纯文本标题, 包含 \cmd{\zthmnumber}, \cmd{\zthmname}, \cmd{\zthmnote} 三
  部分以及一些其它文本. \cs{zthmtitle*} 为 \cs{zthmtitle} 的格式化版本(取消了 CJKecglue, 且可能包含 
  \cs{bfseries}, \cs{sffamily} 等文本格式化命令); 用户在自定义定理类环境样式时应优先使用 \cs{zthmtitle*}, 
  此命令生成的定理类环境标题才能被 \cs{zthmtitleformat} 控制. 此二命令中文本的具体格式可以使用 
  后续的 \cmd{\zthmtitleformat} 命令进行修改.
\end{function}
% \clearpage


\begin{function}[updated=2025-04-25]{\zthmtitleswitch, \zthmtitleswitch*}
  命令 \cmd{\zthmtitleswitch} 用于隐藏定理类环境的标题, 命令 \cmd{\zthmtitleswitch*} 用于显示标题; 在自定义环境
  样式时比较有用. 用户不应该在正文中对此命令进行直接的调用.\par 
  一个基本的使用案例如下\zchcmd:
\end{function}
\vskip1.25em\relax
\begin{DocExample}*
  \begin{theorem}[zthmtitleswitch-1]
    A theorem zthmtitleswitch-1.
  \end{theorem}
  \zthmstylenew{
    ZZZ={begin=, end=, option=\zthmtitleswitch},
  }
  \zthmstyle{ZZZ}
  \begin{theorem}[zthmtitleswitch-2]
    A theorem zthmtitleswitch-2.
  \end{theorem}
\end{DocExample}
关于命令 \cmd{\zthmstyle} 的使用可以参见下面的说明.


\zthmstyle{plain}
\begin{function}[updated=2025-04-25]{\zthmtitleformat, \zthmtitleformat*}
  \begin{syntax}
    \cs{zthmtitleformat}\oarg{type}\marg{format}
  \end{syntax}
  此命令用于修改类型为 \meta{type} 的数学类环境的标题格式(即命令 \cmd{\zthmtitle*} 中的内容), 仅能在导言区使用.%
  \meta{type} 可选值有\texttt{theorem, proof}, 默认值为 \texttt{theorem}.%
  命令 \cs{zthmtitleformat} 仅应用于之后的第一个(类型为 \meta{type} 的)数学类环境标题样式, 而 \cs{zthmtitleformat*} 则应
  用于之后的所有(类型为 \meta{type} 的)数学类环境. \textbf{注意}: 如果 \meta{type} 为 \texttt{proof}, 那么
  在 \meta{format} 中仅有 \cs{zthmname} 和 \cs{thmproof@tmp@color} 可用.\par
  \noindent \textbf{注意}: 此命令默认取消定理类或证明类环境标题中的 CJKecglue, 若用户需要保留这些 glue, 请在格式化代码最前面加上:
  ``\verb|\zthmrestoreCJKecglue|''.
\end{function}
此命令的一个简单使用案例如下\zchcmd:
\ExplSyntaxOn
\tl_set:Nn \l__ztex_proof_color_tl {blue!50}
\ExplSyntaxOff
\begin{DocExample}*
\zthmcolorset{proof=blue!50}
\makeatletter
\zthmtitleformat{\bfseries\color{\thm@tmp@color}\zthmname\zthmnote{\{}{\}}\zthmnumber\ }
\zthmtitleformat[proof]{\bfseries\color{\thmproof@tmp@color}[:\zthmname :]\ }
\makeatother
\begin{theorem}[zthmtitleformat-1]
  A theorem zthmtitleformat-1.
\end{theorem}
\begin{proof}
  This is a proof.
\end{proof}
\end{DocExample}
\ExplSyntaxOn
\tl_set:Nn \l__ztex_proof_color_tl {black}
\ExplSyntaxOff

此外, 还可以参见命令 \cs{zthmnotemptyTF} 中的使用示例. 


\begin{function}[rEXP, updated=2025-04-29]{\zthmnotemptyTF}
  \begin{syntax}
    \cs{zthmnotemptyTF}\marg{true code}\marg{false code}
  \end{syntax}
  此命令用于判断 \cmd{\zthmnote} 是否为空, 如果为空则执行 \meta{true code}, 否则执行 \meta{false code}.
  这个命令在自定义 \cmd{\zthmtitle} 时很有用.\par 
  一个使用样例(\ztex{} 内置的 \texttt{obsidian} 定理样式对应的大致格式, 具体效果可以参见: \cref{sec:theme-library-obsidian}):
\end{function}
\begin{DocExample}
\zthmtitleformat*{\bfseries
  \zthmname\ \zthmnumber
  \zthmnotemptyTF{}{\\}
  \zthmnote{}{}
}
\end{DocExample}



\begin{function}[updated=2025-04-25]{\zthmstylenew}
  \begin{syntax}
    \cs{zthmstylenew}\marg{key-value}
  \end{syntax}
  此命令用于定义新的定理类环境样式, 仅能在导言区使用.
\end{function}


\begin{keyval}[parent=ztex/..]{begin, end, option, preamble}
  \begin{syntax}
    begin     = \meta{code}>\dval{无}
    end       = \meta{code}>\dval{无}
    option    = \meta{code}>\dval{无}
    preamble  = \meta{code}>\dval{无}
  \end{syntax}
  其中 \meta{code} 为任意合法的 \LaTeX{} 代码, 这些代码会被置于对应定理类环境的样式代码中.
  \meta{begin} 和 \meta{end} 即为这个新样式对应环境的开头和结尾; \meta{option} 中的代码
  在 \meta{begin} 之后, 也在环境的开头, 常用于放置一些控制代码; \meta{preamble} 中的代码
  会被 \ztex{} 置于文档的导言区, 常用于放置一些用于定理类环境标题格式化的代码.\par
  当用户声明对应的 \meta{style} 后, 可以在导言区使用命令: \cmd{\zthmstyle}\texttt{\marg{style}} 进行加载.
\end{keyval}
此命令的基本调用格式如下:
\begin{DocExample}[++]
\zthmstylenew
  {
    \ES{\meta{style A}} =
      {
        begin=\ES{\meta{begin code 1}}, 
        end=\ES{\meta{end code 1}}, 
        option=\ES{\meta{option code 1}},
        preamble=\ES{\meta{preamble code 1}}
      },
    \ES{\meta{style B}} =
      {
        begin=\ES{\meta{begin code 2}}, 
        end=\ES{\meta{end code 2}}, 
        option=\ES{\meta{option code 2}},
        preamble=\ES{\meta{preamble code 2}}
      },
    ...
  }
\end{DocExample}


\clearpage
\subsubsection{环境钩子}
\begin{function}[added=2025-05-05]{
    ztex/thm-theorem/before,
    ztex/thm-theorem/begin,
    ztex/thm-theorem/end,
    ztex/thm-theorem/after,
  }
  这四个通用的钩子会将代码添加到所有的 \textbf{定理类} 环境中, 更一般的钩子机制请参见后续
  的 \cmd{\zthmhook} 或 \cmd{\zthmproofhook} 命令.
\end{function}


\begin{function}[added=2025-05-05]{
    ztex/thm-proof/before,
    ztex/thm-proof/begin,
    ztex/thm-proof/end,
    ztex/thm-proof/after,
  }
  这四个通用的钩子会将代码添加到所有的 \textbf{证明类} 环境中, 更一般的钩子机制请参见后续
  的 \cmd{\zthmhook} 或 \cmd{\zthmproofhook} 命令.
\end{function}


\begin{function}[added=2025-05-05]{ztex/thm-theorem/titleformat, ztex/thm-proof/titleformat}
  这两个钩子主要用于修改定理(证明)类环境的标题样式, 前述的 \cmd{\zthmtitleformat} 命令基于这里的
  \texttt{ztex/thm-theorem/titleformat} 钩子.
\end{function}


\begin{function}[updated=2025-04-25]{\zthmhook, \zthmhook*}
  \begin{syntax}
    \cs{zthmhook}\oarg{name}\marg{key-value}
    \cs{zthmhook*}\oarg{name}\marg{key-value}
  \end{syntax}
  此命令用于给已有的(名称为 \meta{name} 的)定理类环境 Hook 中添加代码, \meta{name} 的默认值为 \texttt{theorem}. 
  已有的 Hook: \zkey{ztex/thm/before, ztex/thm/begin, ztex/thm/end, ztex/thm/after}. \cs{zthmhook} 只应用于下
  一个定理类环境, \cs{zthmhook*} 会应用于接下来的所有定理类环境. \texttt{(wrapper begin)} 和 \texttt{(thm-title)}
  相关的钩子请参见 \cs{zthmbefore} 和 \cs{zthmtitlebefore} 命令.  各个 Hook 的位置分布如下:
\end{function}
\def\exampleUR{}
\begin{DocExample}[@@]
(ztex/thm/before) --> (wrapper begin) --> (thm-title) 
  --> (ztex/thm/begin) --> 
      (thm-content) 
  --> (ztex/thm/end) --> 
(wrapper end) --> (ztex/thm/after)
\end{DocExample}
\resetExampleUR

这两个命令不支持手动设置\meta{label}, 针对于 \cmd{\zthmhook*}, \ztex{} 会自动设置 \meta{label}, 其格式
为 \cmd{thm-hook.\meta{Hook Index}}. 


\begin{keyval}[parent=..]{before, begin, end, after}
  \begin{syntax}
    before = \meta{code}>\dval{无}
    begin  = \meta{code}>\dval{无}
    end    = \meta{code}>\dval{无}
    after  = \meta{code}>\dval{无}
  \end{syntax}
  其中 \meta{code} 为合法的 \LaTeX{} 代码片段.
\end{keyval}
一个简单的使用案例如下:
\begin{DocExample}*
\begin{theorem}[zthmhook-1]
  This is a theorem zthmhook-1.
\end{theorem}
\zthmhook{before=ZZa\ , begin=ZZb\ ,}
\begin{theorem}[zthmhook-2]
  This is a theorem zthmhook-2.
\end{theorem}
\end{DocExample}


\begin{function}[updated=2025-04-25]{\zthmproofhook, \zthmproofhook*}
  \begin{syntax}
    \cs{zthmproofhook}\oarg{name}\marg{key-value}
    \cs{zthmproofhook*}\oarg{name}\marg{key-value}
  \end{syntax}
  此命令用于给已有的(名称为 \meta{name} 的)证明类环境 Hook 中添加代码, \meta{name} 的默认值为 \texttt{proof}.
  已有的 Hook: \zkey{ztex/proof/before, ztex/proof/begin, ztex/proof/end, ztex/proof/after}. 
  \cs{zthmproofhook} 只应用于下一个证明类环境, \cs{zthmproofhook*} 会应用于接下来的所有证明类环境.
  各个 Hook 的位置分布如下:
\end{function}
\def\exampleUR{}
\begin{DocExample}[@@]
(ztex/proof/before) --> (proof-title)  
  --> (ztex/proof/begin) --> 
      (proof-content) 
  --> (ztex/proof/end)   -->
(env icon) --> (ztex/proof/after)
\end{DocExample}
\resetExampleUR

和 \cs{zthmhook}, \cs{zthmhook*} 类似, 此二命令会自动设置对应的 \meta{label}, 无需用户手动指定.


\begin{keyval}[parent=..]{before, begin, end, after}
  \begin{syntax}
    before = \meta{code}>\dval{无}
    begin  = \meta{code}>\dval{无}
    end    = \meta{code}>\dval{无}
    after  = \meta{code}>\dval{无}
  \end{syntax}
  其中 \meta{code} 为合法的 \LaTeX{} 代码片段.
\end{keyval}
一个简单的使用样例如下:
\begin{DocExample}*
\zthmproofhook*[solution]{
  before=\noindent\textbf{\color{red}BEFORE},
  begin=\textbf{\color{red}BEGIN},
  end=\textbf{\color{red}END},
  after=\textbf{\color{red}AFTER},
}
\begin{proof}
  This is a proof.
\end{proof}
\begin{solution}
  This is solution I.
\end{solution}
\begin{solution}
  This is solution II.
\end{solution}
\end{DocExample}



\begin{function}[updated=2025-04-25]{\zthmbefore}
  \begin{syntax}
    \cs{zthmbefore}\oarg{type}\marg{code}
  \end{syntax}
  此命令用于把 \meta{code} 置于每个类别为 \meta{type} 的数学环境(如果 \meta{type} 为 \texttt{theorem}, 也就是
  命令 \cmd{\__ztex_thm_warp_start:nnn}; 如果 \meta{type} 为 \texttt{proof}, 那么就是 \cs{__ztex_thm_proof_title:} )之前.
  \meta{type} 的可选值有: \texttt{theorem, proof}, 默认值为 \texttt{theorem}. \meta{code} 默认为 \cs{par}, 用户可以
  把 \meta{code} 置为空, 或设置为 \cs{noindent} 以取消段落缩进.\par 
  一个简单的使用样例如下:
\end{function}
\begin{DocExample}*
\zthmbefore{}
Inline item:%
\begin{theorem}[zthmbefore-1]
  This is a theorem.%
\end{theorem}%
\begin{proposition}[zthmbefore-2]
  This is proposition I.
\end{proposition}
\begin{proof}
  This is a proof.
\end{proof}
\end{DocExample}
\zthmbefore{\par}


\begin{function}[updated=2025-04-25]{\zthmtitlebefore}
  \begin{syntax}
    \cs{zthmtitlebefore}\oarg{type}\marg{code}
  \end{syntax}
  此命令用于把 \meta{code} 置于每个类型为 \meta{type} 的数学环境标题之前. \meta{type} 的可选值有: \texttt{theorem, proof}, 默认
  值为 \texttt{theorem}. \meta{code} 默认为 \cs{noindent}, 用户可以把 \meta{code} 置为空以保留段落缩进.\par 
  一个简单的使用样例如下:
\end{function}
\begin{DocExample}*
\zthmtitlebefore[proof]{[PRF-LIKE]}
\begin{solution}
  This is solution zthmtitlebefore.
\end{solution}
\end{DocExample}
\zthmtitlebefore[proof]{\noindent}
